\documentclass[10pt,a4paper]{amsart}
\usepackage[margin=19mm]{geometry}
\usepackage{amsmath,amssymb,mathtools}
\usepackage[scr=rsfs,cal=euler]{mathalfa}
\usepackage{xfrac}
\usepackage[unicode,hidelinks,bookmarks=false]{hyperref}
\newcommand{\ie}{i.e.\ }
\newcommand{\cf}{cf.\ }
\newcommand{\resp}{respectively\ }
\newcommand{\Subsection}{\subsection}
\providecommand{\nobreakdash}{\nobreak\hbox{-}}
\setlength{\emergencystretch}{2em}
\multlinegap0pt
\usepackage{bm}
\usepackage{bbm}
\usepackage{dsfont}
\usepackage{xspace}
\usepackage{cancel}
\usepackage{mathtools}
\usepackage{amsthm}
\makeatletter
\@ifundefined{thm}{\newtheorem{thm}{Theorem}}{}
\@ifundefined{lem}{\newtheorem{lem}[thm]{Lemma}}{}
\@ifundefined{prop}{\newtheorem{prop}[thm]{Proposition}}{}
\makeatother
\def\R{\mathbb{R}}
\newcommand{\abs}[1]{\lvert #1 \rvert}
\newcommand{\norm}[1]{\lVert #1 \rVert}
\newcommand{\vep}{\varepsilon}
\title[Space-time geometry of Hele-Shaw flow]{Space-time geometry of Hele-Shaw flow}
\author{Carson Collins, Inwon Kim, Sebastian Munoz}
\date{Source: arXiv:2608.16509v1 (CC BY 4.0)}
\begin{document}
\maketitle

\noindent\emph{Source and licence.} Space-time geometry of Hele-Shaw flow. Version arXiv:2608.16509v1 (CC BY 4.0), distributed under the Creative Commons Attribution 4.0 International licence, as recorded in the arXiv licence metadata for this exact version and shown on the abstract page https://arxiv.org/abs/2608.16509v1. Full rights notice: licenses/arxiv-2608.16509-CC-BY-4.0.txt.\par\smallskip

\noindent\textit{Editorial scope.} Counted unit: the Prop whose source label is \texttt{lem:cone\_gradient\_intrinsic\_window} in the article (source0.tex lines 2682--2698, source label \texttt{lem:cone\_gradient\_intrinsic\_window}), delivered together with the article's own complete original proof of it (source0.tex lines 2699--2966, one \texttt{proof} environment of 268 lines). No mathematics is written, repaired or re-derived: statement and proof are the article's own text line for line. Required context retained uncounted: eq:intro-HS-law (lines 78--83); eq:A-source (lines 193--202); eq:A-initial (lines 203--205); eq:A-interior-ball (lines 218--220); eq:A-source-time (lines 226--229); thm:HS-hitting-time-lipschitz (lines 252--277); eq:HS-slice-pressure-comparison (lines 825--828); prop:HS-USC-representative (lines 1091--1110); eq:HS-pointwise-pressure-equation (lines 1628--1631); eq:u\_alpha (lines 1965--1967); eq:cone\_exponent (lines 2096--2098); eq:closing-T-holder (lines 2179--2182); lem:r\_plus\_upper\_bd (lines 2219--2226); eq:cone\_w0\_bounds (lines 2597--2602); eq:cone\_Ur\_positive\_at\_t0 (lines 2605--2608); eq:cone\_f\_slice (lines 2627--2633); lem:cone\_pressure\_away\_spine (lines 2640--2652). Every definition, display and label of those lines is retained unchanged. Omitted from this package and not redistributed: 1--77, 84--192, 206--217, 221--225, 230--251, 278--824, 829--1090, 1111--1627, 1632--1964, 1968--2095, 2099--2178, 2183--2218, 2227--2596, 2603--2604, 2609--2626, 2634--2639, 2653--2681, 2967--3929. Nothing inside a retained excerpt is omitted or altered: every retained body is delivered exactly as the article's lines are written, so no omission receipt of any kind accompanies this package. Citations: the retained excerpts carry no citation command at all, so no bibliography entry of the article is transcribed and no reference is redistributed. Reference adaptation: none was needed and none was made; every reference inside the delivered unit resolves inside it, so no unresolved reference or citation is produced. Printed slips kept literal: no printed word, symbol, hypothesis or number of the counted statement or of its proof was altered. Layout and class: the article's own class is \texttt{amsart} with packages including tikz, inputenc, fontenc, geometry, amsmath, amsthm, amssymb, mathtools, enumitem, xcolor, hyperref; the wrapper is the lane's pinned \texttt{amsart} editorial wrapper at 10pt on A4, which restates the theorem-like environments the retained text needs and adds the article's own macro definitions that the retained text needs (\texttt{\textbackslash R}, \texttt{\textbackslash abs}, \texttt{\textbackslash norm}, \texttt{\textbackslash vep}), with extra wrapper packages (bm, bbm, dsfont, xspace, cancel, mathtools). The delivered numbering is the unit's own. Assets: no retained span contains a figure environment, a graphics-inclusion command, a PDF-page inclusion, a table or a TikZ picture; no image file is copied into this package, the package declares no asset dependency and no image file is redistributed with it. Licence: version arXiv:2608.16509v1 of this article is distributed under the Creative Commons Attribution 4.0 International licence, as recorded in the arXiv licence metadata for this exact version and shown on the abstract page https://arxiv.org/abs/2608.16509v1. A full rights notice is delivered at licenses/arxiv-2608.16509-CC-BY-4.0.txt. Characters: every retained excerpt is pure ASCII, so the delivered engine, XeLaTeX with its default Latin Modern text font, reports no missing character.
\begin{equation*}\tag{HS}\label{eq:intro-HS-law}
\left\{\begin{array}{lll}
    -\Delta p=c &\hbox{ in }&\{p>0\},\\
    V=|\nabla p| &\hbox{on }&\partial\{p>0\},
\end{array}\right.
\end{equation*}

\begin{equation*}\tag{A1}\label{eq:A-source}
\begin{gathered}
    c:\R^d\times[0,\infty)\to[0,\infty)\hbox{ is continuous and, for every }\tau>0,\\
    0<\underline c_\tau\leq c\leq\overline c_\tau<\infty
    \quad\hbox{on }\R^d\times[0,\tau],
    \qquad
    \norm{\nabla_xc(\cdot,t)}_{L^\infty(\R^d)}\leq L_\tau
    \quad\hbox{for a.e. }t\in[0,\tau];
\end{gathered}
\end{equation*}

\begin{equation*}\tag{A2}\label{eq:A-initial}
    \Omega_0^{+}\hbox{ is bounded.}
\end{equation*}

\begin{equation*}\tag{IB}\label{eq:A-interior-ball}
    \Omega_0^{+}\hbox{ satisfies a uniform interior ball condition of radius }r_0>0 .
\end{equation*}

\begin{equation*}\tag{S$_m$}\label{eq:A-source-time}
    (\partial_tc)^-\in L^1_{\operatorname{loc}}\bigl([0,\infty);L^m(\R^d)\bigr),
    \qquad m\in(d/2,\infty] .
\end{equation*}

\begin{thm}[Nondegeneracy and Lipschitz regularity of the hitting time]\label{thm:HS-hitting-time-lipschitz}
Assume \eqref{eq:A-source}, \eqref{eq:A-initial}, \eqref{eq:A-interior-ball}, and \eqref{eq:A-source-time} with $m=\infty$. Then, with $O=T^{-1}((0,\infty)),$
\[
    \partial\{p>0\} =\{(x,T(x)):x\in O\},
\]
and $T:\R^d\to[0,\infty]$ is continuous in the extended sense. Moreover, for each $\tau>0$
there exist $L<\infty$ and $\vep>0$, depending only on
\[
    d,\ r_0,\ \tau,\ \inf_{\R^d\times[0,\tau+1]}c,\
    \norm{\nabla_xc}_{L^\infty(\R^d\times(0,\tau+1))},\
    \norm{(\partial_tc)^-}_{L^1(0,\tau+1;L^\infty(\R^d))},
\]
such that
\begin{equation}\label{eq:HS-T-local-Lipschitz-estimate}
    \abs{T(x)-T(y)}\leq L\abs{x-y}
    \qquad\hbox{whenever }\min\{T(x),T(y)\}<\tau\hbox{ and }\abs{x-y}<\vep .
\end{equation}
In particular, $O$ is open, and $T$ is locally Lipschitz in $O$.
There also exist $\kappa,r_*>0$, depending only on the same quantities and on
$\overline c_\tau:=\sup_{\R^d\times[0,\tau]}c$, such that
\begin{equation}\label{eq:HS-pressure-linear-nondegeneracy}
    \sup_{B_r(x)}p(\cdot,t)\geq\kappa r
    \qquad\hbox{if }x\in\partial\Omega_t,\quad 0<t<\tau,\quad 0<r<r_*.
\end{equation}

\end{thm}

\begin{equation}\label{eq:HS-slice-pressure-comparison}
    p(\cdot,s)\leq Ap(\cdot,t)
    \qquad\hbox{ a.e. in  } \R^d .
\end{equation}

\begin{prop}[Canonical upper semicontinuous representative]\label{prop:HS-USC-representative}
Let \(p\) be a weak solution to \eqref{eq:intro-HS-law}, and assume
\eqref{eq:A-source}, \eqref{eq:A-initial}, and \eqref{eq:A-source-time}.
Then \(p\) has an
upper semicontinuous representative on $Q$, still denoted
by \(p\), characterized by
\begin{equation}\label{eq:HS-forward-average-representative}
    p(x,t)=\lim_{r\downarrow0}
    \frac{1}{r^2\abs{B_r}}
    \int_t^{t+r^2}\int_{B_r(x)}p(y,s)\,dy\,ds  \,\, \hbox{ for } (x,t)\in Q.
\end{equation}
Moreover, for a.e. \(t>0\),
\begin{equation}\label{eq:HS-spatial-envelope-representative}
    p(x,t)=
    \lim_{r\downarrow0}
    \operatorname*{ess\,sup}_{B_r(x)}p(\cdot,t)
    \qquad\hbox{for every }x\in\R^d,
\end{equation}
and \(p\) is right-continuous in time.
\end{prop}

\begin{equation}\label{eq:HS-pointwise-pressure-equation}
    -\Delta p(\cdot,t)=c(\cdot,t)
    \qquad\hbox{in }\mathcal D'(\Omega_t).
\end{equation}

\begin{equation}\label{eq:u_alpha}
    U(V,\alpha,R_0,r_0) := \{ x : r_0 < |x| < R_0 \hbox{ and } |x'| > \alpha |x''| \}.
\end{equation}

\begin{equation}\label{eq:cone_exponent}
    \nu_\alpha := \lambda_\alpha[\tfrac{d-2}{2} + ((\tfrac{d-2}{2})^2 + \lambda_\alpha)^{1/2}]^{-1},
\end{equation}

\begin{equation}\label{eq:closing-T-holder}
    T\in C^{0,\beta}_{\operatorname{loc}}(\{0<T<\infty\})
    \qquad\hbox{for some }\beta\in(0,1] .
\end{equation}

\begin{lem}\label{lem:r_plus_upper_bd}
Assume \eqref{eq:A-source}, \eqref{eq:A-initial}, \eqref{eq:A-source-time}, and
\eqref{eq:closing-T-holder}. Let $R_0=R_0(d,\lambda_1, \alpha,\|c\|_{L^1_t C^{0,\beta}_x}, \|T\|_{C^{0,\beta}})$ be sufficiently small. Then
\begin{equation}\label{eq:upper_bd_r}
    r_+(t) \leq C(t_0 - t)^{1/2},
\end{equation}
where $C = C(\lambda_1, \alpha, \|p\|_{L^\infty(B_{R_0}\times (t,t_0))})$.
\end{lem}

\begin{equation}\label{eq:cone_w0_bounds}
    w(x,t_0)\leq C|x'|^2+o(r^2),
    \qquad
    |\nabla w(x,t_0)|\leq C|x'|+o(r),
    \qquad r\downarrow0.
\end{equation}

\begin{equation}\label{eq:cone_Ur_positive_at_t0}
    w(x,t_0)\geq c_\alpha |x|^2,
    \qquad x\in B_{R_0},\ |x'|\geq \alpha |x''|.
\end{equation}

\begin{equation}\label{eq:cone_f_slice}
    \frac12\leq f(x,s)\leq 1,
    \qquad
    |\nabla_x f(x,s)|\leq L_f,
    \qquad
    (x,s)\in B_{R_0}\times [t_0-R_0^2,t_0+R_0^2].
\end{equation}

\begin{lem}[Pressure lower bound off the spine]\label{lem:cone_pressure_away_spine}
Assume the hypotheses of Theorem~\ref{thm:HS-hitting-time-lipschitz}, and let
\(0\in\Sigma_k(t_0)\) with \(0\leq k\leq d-2\).
Write \(U(\alpha,R_0,r):=U(V,\alpha,R_0,r)\), as in \eqref{eq:u_alpha}. Set \(\gamma=\nu_\alpha\), with \(\nu_\alpha\) defined in \eqref{eq:cone_exponent}, if \(k\geq1\), and fix \(\gamma\in(0,1)\) if \(k=0\). Then for \(\delta\in(4\alpha,1)\), there exists $c_*=c_*(d,k,\alpha,\delta,\gamma,R_0, \underline c)>0$ with the following
property:  for any $r\in (0,\frac{R_0}{128})$, if 
\[
    U(\alpha,R_0,r)\subset\Omega_s \quad \hbox{ for some  }
s\in [t_0-R_0^2, t_0+R_0^2], \]
then, for any $R\in [16r, 32r]$, 
\begin{equation}\label{eq:cone_pressure_lower_away_spine}
    p(x,s)\geq c_*R^\gamma \hbox{ in } \quad \{R/2\leq |x|\leq R\}\cap \{|x'| \geq \delta R\}.
\end{equation}
\end{lem}

\begin{prop}[Intrinsic window gradient estimate]\label{lem:cone_gradient_intrinsic_window}
Assume the hypotheses of Theorem~\ref{thm:HS-hitting-time-lipschitz}, let
\(0\in\Sigma_k(t_0)\) with \(0\leq k\leq d-2\), and let \(0<R_0<R_*\) be as fixed above.
There exists
\(\bar\alpha\in(0,1)\), depending only on \(d,k\), the constant in
\eqref{eq:cone_w0_bounds}, and the constants in \eqref{eq:cone_f_slice}, such
that if \(0<\alpha\leq\bar\alpha\), then the following holds.
Let \(\gamma=\nu_\alpha\) if \(k\geq1\), and let
\(\gamma\in(0,1)\) if \(k=0\).
For every \(\sigma\in(\gamma,1)\), there exist constants
\(c_\sigma\in(0,1)\) and \(r_\sigma\in(0,R_0)\) such that, for every \(R\in(0,r_\sigma)\),
\begin{equation}\label{eq:cone_gradient_intrinsic_window}
    |\nabla w(x,t)|\leq R^{1-\sigma} p(x,t),
    \qquad
    (x,t)\in B_{R/4}\times [t_0-c_\sigma R^2,t_0+c_\sigma R^{2-\sigma}].
\end{equation}
\end{prop}

\begin{proof}
Fix \(\sigma\in(\gamma,1)\) and fix a unit vector \(e\). It suffices to estimate the directional derivative $\partial_e w$, with constants independent of \(e\). We begin by choosing some parameters. Let
\[
    K:=
    \left(\frac{6}{\lambda_1\alpha^2}
    \|p\|_{L^\infty(B_{R_0}\times(t_0-R_0^2,t_0))}
    \right)^{1/2}.
\]
Choose \(c_\sigma\in(0,1)\) so small that
\begin{equation}\label{eq:cone_csigma_choice}
    K\sqrt{c_\sigma}\leq \frac1{32}.
\end{equation} Next choose \(r_\sigma\in(0,R_0/4)\) so small that, whenever
\(0<R<r_\sigma\),
\[
    [t_0-c_\sigma R^2,t_0+c_\sigma R^{2-\sigma}]
    \subset (t_0-R_0^2,t_0+R_0^2).
\]
Thus \eqref{eq:cone_f_slice} applies for all times in this interval.


Fix \(x_0\in B_{R/4}\) and \(t\in[t_0-c_\sigma R^2,t_0+c_\sigma R^{2-\sigma}]\). If
\(w(x_0,t)=0\), then \(\nabla w(x_0,t)=0\), and there is nothing to prove. We therefore assume
\(w(x_0,t)>0\). Let \(\varepsilon_j\downarrow0\), with
\(\varepsilon_j<w(x_0,t)\), and choose smooth open sets $\Omega_j$ such that 
\begin{equation}\label{eq:cone_Omega_j}
    B_{R/2}\cap\{w(\cdot,t)>\varepsilon_j\}\subset\Omega_j
    \subset B_R\cap\Omega_t.
\end{equation}

For nonnegative \(C^{1,1}\) functions we use
\begin{equation}\label{eq:cone_c11_sqrt}
    |\nabla u|^2\leq C u,
\end{equation}
with \(C\) controlled by the local \(C^{1,1}\) norm. Applying this to
\(u=w(\cdot,t)\) gives
\begin{equation}\label{eq:cone_inner_boundary_smallness}
    0\leq w(x,t)\leq\varepsilon_j,
    \qquad
    |\partial_e w(x,t)|\leq C\sqrt{\varepsilon_j},
    \qquad x\in\partial\Omega_j\cap B_{R/2}.
\end{equation}

Let \(G_{\Omega_j}\) be the Green function for \(-\Delta\) in \(\Omega_j\), and
set
\[
    I_j:=\int_{\Omega_j}G_{\Omega_j}(x_0,y)\,dy,
\]
\[
    c_j:=
    \frac{\int_{\Omega_j}f(y,t)G_{\Omega_j}(x_0,y)\,dy}{I_j},
    \qquad
    a_j:=
    \frac{\int_{\Omega_j}\partial_e f(y,t)G_{\Omega_j}(x_0,y)\,dy}{I_j}.
\]
We define the correctors \(\psi_j\) and \(\psi_j^e\) by
\[
    \Delta\psi_j=f(\cdot,t)-c_j,
    \qquad
    \Delta\psi_j^e=\partial_e f(\cdot,t)-a_j,
    \qquad
    \psi_j=\psi_j^e=0\hbox{ on }\partial\Omega_j.
\]
By \eqref{eq:cone_f_slice}, \(c_j\in[1/2,1]\) and \(|a_j|\leq L_f\).  The
normalization by the Green function gives
\begin{equation}\label{eq:cone_corrector_center}
    \psi_j(x_0)=\psi_j^e(x_0)=0.
\end{equation}

Define
\[
    Q(x):=\frac{|x-x_0|^2}{2d},
    \qquad
    q_j:=\partial_e w(\cdot,t)-\psi_j^e-a_jQ,
\]

With \(\kappa>0\) to be chosen,
we consider the following function:

\begin{equation}\label{eq:cone_barrier}
    z_j:=p(\cdot,t)-R^{-(1-\sigma)}q_j
    +\kappa R^{-(2-\sigma)}\bigl(c_jQ-w(\cdot,t)+\psi_j\bigr).
\end{equation}
Note that $z_j$ is superharmonic in $\Omega_j$, since  \(\Delta w(\cdot,t)=f(\cdot,t)\) and
the pressure equation \eqref{eq:HS-pointwise-pressure-equation}
give
\[
    \Delta q_j=0,
    \qquad
    \Delta\bigl(c_jQ-w(\cdot,t)+\psi_j\bigr)=0 \qquad \hbox{ in } \Omega_j.
\]
Our goal is to estimate the lower bound of $z_j$ on $\partial\Omega_j$, which then yields the desired estimate at $x=x_0$ by superharmonicity.
On \(\partial\Omega_j\), the correctors vanish, and therefore
\[
    z_j=p(\cdot,t)-R^{-(1-\sigma)}\partial_e w(\cdot,t)
    +\bigl(\kappa c_jR^{-(2-\sigma)}+a_jR^{-(1-\sigma)}\bigr)Q
    -\kappa R^{-(2-\sigma)}w(\cdot,t).
\]
After \(\kappa\) is fixed, we decrease \(r_\sigma\) so that
\(|a_j|R\leq\kappa/4\) for \(R<r_\sigma\). Since \(c_j\geq1/2\),
\begin{equation}\label{eq:cone_boundary_common}
    z_j\geq p(\cdot,t)-R^{-(1-\sigma)}\partial_e w(\cdot,t)
    +\frac{\kappa}{4}R^{-(2-\sigma)}Q
    -\kappa R^{-(2-\sigma)}w(\cdot,t)
\end{equation}
on \(\partial\Omega_j\). On \(\partial\Omega_j\cap B_{R/2}\), this and
\eqref{eq:cone_inner_boundary_smallness} give
\begin{equation}\label{eq:cone_inner_boundary_zj}
    z_j\geq
    -CR^{-(1-\sigma)}\sqrt{\varepsilon_j}
    -\kappa R^{-(2-\sigma)}\varepsilon_j.
\end{equation}

It remains to estimate \(z_j\) on the outer boundary 
\(\partial\Omega_j\cap\{|x|\geq R/2\}\), where $Q\geq c_dR^2$. For this part we proceed by different arguments at past times $t<t_0$, using results from the barrier argument, and at future times $t>t_0$, using the right continuity of $p$. 


\textbf{\emph{Backward window.}}
Let $x\in \partial\Omega_j\cap \{|x|\geq R/2\}$.
Lemma~\ref{lem:r_plus_upper_bd} gives
\begin{equation}\label{eq:cone_r_plus_small_on_window}
    r_+(t)\leq K\sqrt{t_0-t}\leq \frac{R}{32},
    \qquad t\in[t_0-c_\sigma R^2,t_0].
\end{equation}

 Define
\[
    v:=w(\cdot,t_0)-w(\cdot,t)\geq0.
\]
Using \eqref{eq:cone_c11_sqrt} for \(v\), followed by Young's inequality, we have
\begin{equation}\label{eq:cone_backward_absorb}
    \kappa R^{-(2-\sigma)}v+R^{-(1-\sigma)}\partial_ev
    \geq -C\kappa^{-1}R^\sigma.
\end{equation}
Combining \eqref{eq:cone_boundary_common}, \(w(\cdot,t)=w(\cdot,t_0)-v\),
\eqref{eq:cone_w0_bounds}, and \eqref{eq:cone_backward_absorb}, we obtain
\begin{equation}\label{eq:cone_backward_outer_common}
    \begin{aligned}
    z_j(x)\geq{}&p(x,t)+c_d\kappa R^\sigma
    -CR^{-(1-\sigma)}|x'|
    -C\kappa R^{-(2-\sigma)}|x'|^2\\
    &-o(R^\sigma)-\kappa o(R^\sigma)
    -C\kappa^{-1}R^\sigma-CR^{1+\sigma}.
    \end{aligned}
\end{equation}

We now choose the threshold \(\delta_0\) separating the near-spine and
away-from-spine parts of the boundary in our case analysis. Let \(C_*\) bound
the coefficient of \(\kappa R^{-(2-\sigma)}|x'|^2\) in
\eqref{eq:cone_backward_outer_common}. Take \(\bar\alpha\) small enough that \(8\bar\alpha<1\) and
\(C_*(8\bar\alpha)^2\leq c_d/4\), and choose
\(\delta_0\in(4\alpha,8\bar\alpha)\). Then
\begin{equation}\label{eq:cone_delta_choice}
    4\alpha<\delta_0,
    \qquad
    C_*\delta_0^2\leq c_d/4 .
\end{equation}
Let \(c_*\) be the constant from
Lemma~\ref{lem:cone_pressure_away_spine} with \(\delta=\delta_0\). If \(|x'|\leq\delta_0R\), then
\begin{equation}\label{eq:cone_near_backward}
    z_j(x)\geq p(x,t)
    +\bigl(c_d\kappa-C\delta_0-C\kappa\delta_0^2-C\kappa^{-1}-o(1)-\kappa o(1)\bigr)R^\sigma
    -CR^{1+\sigma}.
\end{equation}
By \eqref{eq:cone_delta_choice}, the coefficient of \(R^\sigma\) in
\eqref{eq:cone_near_backward} is at least
\[
    \frac{c_d}{2}\kappa-C\delta_0-C\kappa^{-1}-o(1)-\kappa o(1),
\]
which is positive once \(\kappa\) is chosen large and \(r_\sigma\) is
decreased so that the \(o(1)\) terms are small for \(R<r_\sigma\).





If \(|x'|\geq\delta_0R\), set
\(r:=r_+(t)+R/32\). By
\eqref{eq:cone_r_plus_small_on_window}, \(R/32\leq r\leq R/16\), and by
the definition of \(r_+\), \(U(\alpha,R_0,r)\subset\Omega_t\).
Lemma~\ref{lem:cone_pressure_away_spine}, applied with \(s=t\), gives
\[
    p(x,t)\geq c_*R^\gamma.
\]
Since \(|x'|\leq R\), \eqref{eq:cone_backward_outer_common} yields
\begin{equation}\label{eq:cone_far_backward}
    z_j(x)\geq
    c_*R^\gamma-C(\kappa+\kappa^{-1})R^\sigma-CR^{1+\sigma}.
\end{equation}
Because \(\gamma<\sigma\), \eqref{eq:cone_near_backward} and
\eqref{eq:cone_far_backward} are nonnegative once \(r_\sigma\) is small enough.

The minimum principle and \eqref{eq:cone_inner_boundary_zj} give
\[
    z_j(x_0)\geq
    -CR^{-(1-\sigma)}\sqrt{\varepsilon_j}
    -\kappa R^{-(2-\sigma)}\varepsilon_j.
\]
By \eqref{eq:cone_corrector_center},
\[
    z_j(x_0)=p(x_0,t)-R^{-(1-\sigma)}\partial_e w(x_0,t)
    -\kappa R^{-(2-\sigma)}w(x_0,t).
\]
Letting \(j\to\infty\) gives
\[
    \partial_e w(x_0,t)\leq R^{1-\sigma}p(x_0,t).
\]
Replacing \(e\) by \(-e\), we conclude \eqref{eq:cone_gradient_intrinsic_window} for $t\leq t_0$.


\textbf{\emph{Forward window.}}
We next estimate $z_j$ on $\partial\Omega_j \cap \{|x|\geq R/2\}$ for the forward times $t>t_0$ with the same choice of $\delta_0$ and $\kappa$ as for backward times, but with potentially smaller choices of \(c_\sigma\) and
\(r_\sigma\) made below. 
If \(t=t_0+\tau\) with
\(\tau\leq c_\sigma R^{2-\sigma}\), set
\[
    v:=w(\cdot,t)-w(\cdot,t_0)=\int_{t_0}^t p(\cdot,s)\,ds\geq0.
\]
Integrating the slice comparison
\eqref{eq:HS-slice-pressure-comparison} in \(s\in(t_0,t)\), and passing to the
canonical representative from Proposition~\ref{prop:HS-USC-representative}, gives
\begin{equation}\label{eq:cone_forward_v_bound}
    v(x)=\int_{t_0}^t p(x,s)\,ds
    \leq \frac{\overline c}{\underline c}\tau p(x,t).
\end{equation}
By \eqref{eq:cone_forward_v_bound} and \(\tau\leq c_\sigma R^{2-\sigma}\), we have
\(v\leq(\overline c/\underline c)\,c_\sigma R^{2-\sigma}p(\cdot,t)\). We decrease
\(c_\sigma\), depending on \(\kappa\) and \(\overline c/\underline c\), so that
\(C\kappa\,(\overline c/\underline c)\,c_\sigma\leq\tfrac12\), where \(C\) is the constant
in \eqref{eq:cone_forward_absorb} below.
Since the \(C^{1,1}\) estimate \eqref{eq:cone_c11_sqrt} for \(v\) gives
\begin{equation}\label{eq:cone_forward_absorb}
    -\kappa R^{-(2-\sigma)}v-R^{-(1-\sigma)}\partial_ev
    \geq
    -C\kappa^{-1}R^\sigma-C\kappa R^{-(2-\sigma)}v,
\end{equation}
it follows 
\begin{equation}\label{eq:cone_forward_outer_common}
    \begin{aligned}
    z_j(x)\geq{}&\frac12p(x,t)+c_d\kappa R^\sigma
    -CR^{-(1-\sigma)}|x'|
    -C\kappa R^{-(2-\sigma)}|x'|^2\\
    &-o(R^\sigma)-\kappa o(R^\sigma)
    -C\kappa^{-1}R^\sigma-CR^{1+\sigma}.
    \end{aligned}
\end{equation}

When \(|x'|\leq\delta_0R\), the choices of \(\delta_0\) in
\eqref{eq:cone_delta_choice}, of \(\kappa\) after
\eqref{eq:cone_near_backward}, and of \(c_\sigma\)  make the right-hand side of
\eqref{eq:cone_forward_outer_common} nonnegative, after the final decrease of the range of $R$ given by 
\(r_\sigma\). When \(|x'|\geq\delta_0R\), set \(h=R/20\). By
\eqref{eq:cone_Ur_positive_at_t0} and monotonicity,
\[
    U(\alpha,R_0,h)\subset\Omega_{t_0}\subset\Omega_t,
\]
and Lemma~\ref{lem:cone_pressure_away_spine}, applied with \(s=t\) and \(r=h\),
gives
\[
    p(x,t)\geq c_*R^\gamma.
\]
Since \(\gamma<\sigma\), the positive pressure term in
\eqref{eq:cone_forward_outer_common} dominates the remaining \(R^\sigma\) and
\(R^{1+\sigma}\) errors for \(R<r_\sigma\). From here the minimum
principle and the passage \(j\to\infty\) proceed exactly as in the backward
window, yielding \(\partial_e w(x_0,t)\leq R^{1-\sigma}p(x_0,t)\); replacing \(e\)
by \(-e\) and taking the supremum over unit directions proves
\eqref{eq:cone_gradient_intrinsic_window}.
\end{proof}

\end{document}
